% Copyright (c) 2008-2009 solvethis
% Copyright (c) 2010-2016,2018-2019 Casper Ti. Vector
% Public domain.
%
% 使用前请先仔细阅读 pkuthss 和 biblatex-caspervector 的文档，
% 特别是其中的 FAQ 部分和用红色强调的部分。
% 两者可在终端/命令提示符中用
%   texdoc pkuthss
%   texdoc biblatex-caspervector
% 调出。

% 采用了自定义的（包括大小写不同于原文件的）字体文件名，
% 并改动 ctex.cfg 等配置文件的用户请自行加入 nofonts 选项；
% 其它用户不用加入 nofonts 选项，加入之后反而会产生错误。
\documentclass[UTF8]{pkuthss}
% 如果的确须要使脚注按页编号的话，可以去掉后面 footmisc 包的注释。
% 注意：在启用此设定的情况下，可能要多编译一次以产生正确的脚注编号。
%\usepackage[perpage]{footmisc}

% 使用 biblatex 排版参考文献，并规定其格式（详见 biblatex-caspervector 的文档）。
% （如果须要使用作者–年编码制，应将 caspervector 换 成 caspervector-ay）
% 这里按照西文文献在前，中文文献在后排序（“sorting = ecnyt”）；
% 若须按照中文文献在前，西文文献在后排序，请设置“sorting = cenyt”；
% 若须按照引用顺序排序，请设置“sorting = none”。
% 若须在排序中实现更复杂的需求，请参考 biblatex-caspervector 的文档。
% \usepackage[backend = biber, style = caspervector, utf8, sorting = none, seconds = true]{biblatex}
\usepackage[backend = biber, style = gb7714-2015, gbnamefmt=quanpin, gbpub=false, doi=false, url=false,eprint=false,gbpunctin=false]{biblatex}
% gbpunctin 控制是否输出两个反斜杠 //
% gbnamefmt 控制人名如何拼写

\usepackage{multirow}
\usepackage{bbm}
\usepackage{tabularx}

\usepackage[justification=centering]{caption} % caption多行时中心对齐失效，用这个可以解决
\usepackage[list=off]{bicaption} % 双语caption
\captionsetup[figure][bi-second]{name=Figure} % 设置图的英文编号前缀
\captionsetup[table][bi-second]{name=Table} % 设置表的英文编号前缀
\usepackage{enumitem} % 科研成果列表
\usepackage{etoolbox}
\AtBeginEnvironment{tabular}{\small} % 默认表格内容的字体有点大，改成是11pt
\special{dvipdfmx:config z 0} % 取消PDF压缩，加快编译速度; 最终版记得取消，不然生成文件会较大，虽然也没什么大不了的。


\renewcommand{\listfigurename}{图目录}
\renewcommand{\listtablename}{表目录}

% 对于 linespread 值的计算过程有兴趣的同学可以参考 pkuthss.cls。
\renewcommand*{\bibfont}{\zihao{5}\linespread{1.27}\selectfont}
% 按学校要求设定参考文献列表的段间距。
\setlength{\bibitemsep}{3bp}

% 设定文档的基本信息。
\pkuthssinfo{
	cthesisname = {硕士研究生学位论文}, ethesisname = {Master Thesis},
	ctitle = {基于PKUTHSS的测试文档}, etitle = {Test Document},
	cauthor = {作者名称},
	eauthor = {Author Name},
	studentid = {19012345678},
	date = {二〇二〇年三月},
	school = {学院名称},
	cmajor = {专业名称}, emajor = {Major Name},
	direction = {方向名称},
	cmentor = {教授名称}, ementor = {Prof.\ Someone},
	ckeywords = {其一，其二}, ekeywords = {First, Second}
}
% 载入参考文献数据库（注意不要省略“.bib”）。
\addbibresource{thesis.bib}

% 普通用户可删除此段，并相应地删除 chap/*.tex 中的
% “\pkuthssffaq % 中文测试文字。”一行。
\usepackage{color}
\def\pkuthssffaq{%
	\emph{\textcolor{red}{pkuthss 文档模版最常见问题：}}

	\texttt{\string\cite}、\texttt{\string\parencite} %
	和 \texttt{\string\supercite} 三个命令分别产生%
	上标且带方括号、带方括号的和未格式化的引用标记：%
	\cite{test-en}，\parencite{test-zh}、\supercite{test-en, test-zh}。

	若要避免章末空白页，请在调用 pkuthss 文档类时加入 \texttt{openany} 选项。

	如果编译时不出参考文献，
	请参考 \texttt{texdoc pkuthss}“问题及其解决”一章
	“上游宏包可能引起的问题”一节中关于 biber 的说明。%
}

\begin{document}
	% 以下为正文之前的部分，默认不进行章节编号。
	\frontmatter
	% 此后到下一 \pagestyle 命令之前不排版页眉或页脚。
	\pagestyle{empty}
	% 自动生成封面。
	\maketitle
	% 版权声明。封面要求单面打印，故须新开右页。
	\cleardoublepage
	\include{chap/copy}

	% 此后到下一 \pagestyle 命令之前正常排版页眉和页脚。
	\cleardoublepage
	\pagestyle{plain}
	% 重置页码计数器，用大写罗马数字排版此部分页码。
	\setcounter{page}{0}
	\pagenumbering{Roman}

	% group 的作用是将\let\cleardoublepage\relax的效果限制在正文内。
	\begingroup

	\let\cleardoublepage\relax

	% 中西文摘要。
	% Copyright (c) 2014,2016 Casper Ti. Vector
% Public domain.

\begin{cabstract}
	\pkuthssffaq % 中文测试文字。
	需要注意的是摘要不可以引用文献。
\end{cabstract}


\clearpage

\begin{eabstract}
	Test of the English abstract.
\end{eabstract}

% vim:ts=4:sw=4

	% 自动生成目录。
	\tableofcontents

	\clearpage
	\listoffigures % 图目录
	\clearpage
	\listoftables % 表目录
	
	\hiddenchap{缩写表}

\begin{tabularx}{\linewidth}{{>{\hsize=.4\hsize}X>{\hsize=1.5\hsize}X>{\hsize=1.1\hsize}X}}
    AI & Artificial Intelligence & 人工智能 \\
\end{tabularx}
 % 术语表 / 缩写表


	% 以下为正文部分，默认要进行章节编号。
	\mainmatter

	% 各章节。
	% Copyright (c) 2014,2016,2018 Casper Ti. Vector
% Public domain.

\chapter{引言}

This is an example sentence.

\section{引用}
这篇工作\cite{vaswani2017attention}不过如此。
\footnote{此处引用的数字没有出错，因为第一个引用在摘要。}

\section{图表}

\subsection{表}

\begin{table}[htbp]
    \centering
    \bicaption{测试表格}{Test table}
    \label{tab:example}
    \begin{tabular}{ll}
    \toprule
    \textbf{写} & \textbf{几句} \\
    \midrule

    表格 & 就  \\
    
    长  & 这样 \\
    
    \bottomrule

    \end{tabular}
\end{table}

\subsection{图}

\begin{figure}[htb]
    \centering
        \includegraphics[width=0.4\textwidth]{image/example.jpeg}

    % \caption{ParAhLSTMat的解码器架构}
    \bicaption{图片示例}{Example of graphics}
    \label{fig:2-par_arch}
\end{figure}

\subsection{公式}

\begin{equation}
    E = mc^{2}
\end{equation}

\pkuthssffaq % 中文测试文字。

% vim:ts=4:sw=4
	\include{chap/chap2}
	\include{chap/chap3}

	% 正文中的附录部分。
	\appendix
	% 排版参考文献列表。bibintoc 选项使“参考文献”出现在目录中；
	% 如果同时要使参考文献列表参与章节编号，可将“bibintoc”改为“bibnumbered”。
	\printbibliography[heading = bibintoc]
	
	% 各附录。
	% \include{chap/encl1}

	\endgroup

	% 以下为正文之后的部分，默认不进行章节编号。
	\backmatter
	
	% 科研成果。
	\chapter{攻读硕士学位期间的科研成果}

\section*{发表论文}

% 此处格式应当和参考文献的格式一致。
\begin{enumerate}[label={[\arabic*]}]
    \item \underline{\textbf{Vaswani Ashish}}, Shazeer Noam, Parmar Niki, et al. Attention is all you need[J]. Advances in neural information processing systems, 2017, 30.
  \end{enumerate}



\section*{参与项目}

\begin{enumerate}[label={[\arabic*]}]
  \item 国家自然科学基金项目——“如何一顿吃饱全家不饿”，2020.01.01--2022.06.01

\end{enumerate}

% \section*{专利申请}


	\clearpage{\thispagestyle{empty}\cleardoublepage}
	
	% 致谢。
	\include{chap/ack}
	\clearpage{\thispagestyle{empty}\cleardoublepage}

	% 原创性声明和使用授权说明。
	% Copyright (c) 2008-2009 solvethis
% Copyright (c) 2010-2017 Casper Ti. Vector
% All rights reserved.
%
% Redistribution and use in source and binary forms, with or without
% modification, are permitted provided that the following conditions are
% met:
%
% * Redistributions of source code must retain the above copyright notice,
%   this list of conditions and the following disclaimer.
% * Redistributions in binary form must reproduce the above copyright
%   notice, this list of conditions and the following disclaimer in the
%   documentation and/or other materials provided with the distribution.
% * Neither the name of Peking University nor the names of its contributors
%   may be used to endorse or promote products derived from this software
%   without specific prior written permission.
%
% THIS SOFTWARE IS PROVIDED BY THE COPYRIGHT HOLDERS AND CONTRIBUTORS "AS
% IS" AND ANY EXPRESS OR IMPLIED WARRANTIES, INCLUDING, BUT NOT LIMITED TO,
% THE IMPLIED WARRANTIES OF MERCHANTABILITY AND FITNESS FOR A PARTICULAR
% PURPOSE ARE DISCLAIMED. IN NO EVENT SHALL THE COPYRIGHT HOLDER OR
% CONTRIBUTORS BE LIABLE FOR ANY DIRECT, INDIRECT, INCIDENTAL, SPECIAL,
% EXEMPLARY, OR CONSEQUENTIAL DAMAGES (INCLUDING, BUT NOT LIMITED TO,
% PROCUREMENT OF SUBSTITUTE GOODS OR SERVICES; LOSS OF USE, DATA, OR
% PROFITS; OR BUSINESS INTERRUPTION) HOWEVER CAUSED AND ON ANY THEORY OF
% LIABILITY, WHETHER IN CONTRACT, STRICT LIABILITY, OR TORT (INCLUDING
% NEGLIGENCE OR OTHERWISE) ARISING IN ANY WAY OUT OF THE USE OF THIS
% SOFTWARE, EVEN IF ADVISED OF THE POSSIBILITY OF SUCH DAMAGE.

{
	\ctexset{section = {
		format+ = {\centering}, beforeskip = {40bp}, afterskip = {15bp}
	}}

	% 学校书面要求本页面不要页码，但在给出的 Word 模版中又有页码且编入了目录。
	% 此处以 Word 模版为实际标准进行设定。
	\chapter{北京大学学位论文原创性声明和使用授权说明}

	\thispagestyle{empty} % 本页无页码和页头

	\vspace*{1em}
	\section*{原创性声明}
	\vspace*{1em}

	本人郑重声明：
	所呈交的学位论文，是本人在导师的指导下，独立进行研究工作所取得的成果。
	除文中已经注明引用的内容外，
	本论文不含任何其他个人或集体已经发表或撰写过的作品或成果。
	对本文的研究做出重要贡献的个人和集体，均已在文中以明确方式标明。
	本声明的法律结果由本人承担。
	\vskip 1em
	\rightline{%
		论文作者签名：\hspace{5em}%
		日期：\hspace{2em}年\hspace{2em}月\hspace{2em}日%
	}

	\section*{%
		学位论文使用授权说明\\[-0.33em]
		\textmd{\zihao{5}（必须装订在提交学校图书馆的印刷本）}%
	}

	本人完全了解北京大学关于收集、保存、使用学位论文的规定，即：
	\begin{itemize}
		\item 按照学校要求提交学位论文的印刷本和电子版本；
		\item 学校有权保存学位论文的印刷本和电子版，
			并提供目录检索与阅览服务，在校园网上提供服务；
		\item 学校可以采用影印、缩印、数字化或其它复制手段保存论文；
		\item 因某种特殊原因须要延迟发布学位论文电子版，
			授权学校$\square$\nobreakspace{}一年/
			$\square$\nobreakspace{}两年 /
			$\square$\nobreakspace{}三年以后在校园网上全文发布。
	\end{itemize}
	\vspace*{\baselineskip}
	\centerline{（保密论文在解密后遵守此规定）}
	\vskip 3em
	\centerline{%
		论文作者签名：\hspace{5em}导师签名：\hspace{5em}%
	}
	\begin{center}
		日期：\hspace{2em}年\hspace{2em}月\hspace{2em}日%
	\end{center}

	% 若须排版二维码，请将二维码图片重命名为“barcode”，
	% 转为合适的图片格式，并放在当前目录下，然后去掉下面 2 行的注释。
	% \vfill\noindent
	% \includegraphics[height = 3em]{barcode}
}

% vim:ts=4:sw=4

	
\end{document}

% vim:ts=4:sw=4